% ── §2.6: Αξιολόγηση Συνεισφοράς (Shapley Values) ───────────

\section{Αξιολόγηση Συνεισφοράς (Shapley Values)}
\label{sec:shapley}

Ένα κεντρικό ερώτημα στα συνεργατικά συστήματα είναι: \textit{πόσο
συνέβαλε κάθε συμμετέχων στο συνολικό αποτέλεσμα;} Στο πλαίσιο του
FL, αυτό αντιστοιχεί στο ερώτημα πόσο βελτίωσε κάθε client το
καθολικό μοντέλο. Η απάντηση δίνεται μέσω της θεωρίας Συνεργατικών
Παιγνίων (Cooperative Game Theory) και ειδικότερα μέσω των
\textit{τιμών Shapley} \cite{Shapley1953}.

\subsection{Συνεργατικά Παίγνια και Συναρτήσεις Χαρακτηριστικής}
\label{subsec:cooperative_games}

Ένα συνεργατικό παίγνιο ορίζεται ως ζεύγος $(N, v)$, όπου
$N = \{1, \ldots, n\}$ το σύνολο παικτών (clients) και
$v: 2^N \rightarrow \mathbb{R}$ η \textbf{συνάρτηση χαρακτηριστικής}
(characteristic function) που αποδίδει μια αξία σε κάθε υποσύνολο
$S \subseteq N$ \cite{Wang2019}. Στο HAR-FL σύστημά μας, $v(S)$
αντιστοιχεί στην ακρίβεια του μοντέλου που θα προέκυπτε αν μόνο
το υποσύνολο $S$ clients συμμετείχε στην εκπαίδευση.

\subsection{Ορισμός και Ιδιότητες Shapley Value}
\label{subsec:shapley_def}

Η τιμή Shapley $\phi_i(v)$ αποδίδει στον παίκτη $i$ το σταθμισμένο
μέσο περιθώριο συνεισφοράς σε όλες τις δυνατές σειρές εισόδου:

\begin{equation}
    \label{eq:shapley}
    \phi_i(v) = \sum_{S \subseteq N \setminus \{i\}}
    \frac{|S|!\,(|N|-|S|-1)!}{|N|!}
    \bigl[v(S \cup \{i\}) - v(S)\bigr]
\end{equation}

\noindent Ο παράγοντας $[v(S \cup \{i\}) - v(S)]$ μετράει την
\textit{οριακή συνεισφορά} (marginal contribution) του client $i$
όταν εντάσσεται στο υποσύνολο $S$. Ο συντελεστής
$|S|!\,(|N|-|S|-1)!/|N|!$ σταθμίζει αναλόγως με την πιθανότητα
εμφάνισης αυτής της σειράς.

Οι τιμές Shapley ικανοποιούν τέσσερις αξιωματικές ιδιότητες:
\textit{αποδοτικότητα} (το σύνολο της αξίας κατανέμεται πλήρως),
\textit{συμμετρία} (ίδια συνεισφορά → ίδια ανταμοιβή),
\textit{μηδενικός παίκτης} (μηδενική συνεισφορά → μηδενική ανταμοιβή)
και \textit{προσθετικότητα}. Αυτές οι ιδιότητες εγγυώνται τη
\textbf{δικαιοσύνη} (fairness) της κατανομής ανταμοιβών \cite{Shapley1953}.

\subsection{Υπολογιστική Πολυπλοκότητα και Προσεγγίσεις}
\label{subsec:shapley_approx}

Ο ακριβής υπολογισμός της Eq.~\eqref{eq:shapley} απαιτεί
αξιολόγηση $v(S)$ για $2^n$ υποσύνολα — εκθετική πολυπλοκότητα
που καθιστά τον ακριβή υπολογισμό απαγορευτικό για $n > 20$.
Στην πράξη χρησιμοποιείται η \textbf{προσέγγιση Leave-One-Out (LOO)}:

\begin{equation}
    \label{eq:loo}
    \hat{\phi}_i = v(N) - v(N \setminus \{i\})
\end{equation}

\noindent η οποία υπολογίζει την οριακή συνεισφορά κάθε client
ως διαφορά ακρίβειας με και χωρίς τον client $i$, απαιτώντας μόνο
$n+1$ αξιολογήσεις \cite{Gopinathan2023}. Εναλλακτικά, η
\textit{gradient-based} προσέγγιση εκτιμά τη συνεισφορά από τις
κλίσεις (gradients) χωρίς επαναξιολόγηση μοντέλου \cite{Pavlidis2021}.

Στην παρούσα εργασία, η LOO προσέγγιση χρησιμοποιείται ως η
ποσότητα $q_i$ του μηχανισμού κινήτρων (§~\ref{sec:game_theory}),
υπολογιζόμενη στον server μετά από κάθε γύρο FL
(βλ.\ Ενότητα~\ref{sec:incentive_mechanism}).
